\documentclass[aspectratio=169,10pt]{beamer}
% ═══════════════════════════════════════════════════════════════════════════════
% SHARED PREAMBLE — Foundation Models Course (HIT)
% To be \input by each lecture file AFTER \documentclass
% ═══════════════════════════════════════════════════════════════════════════════

% ─────────────────────────────────────────────────────────────────────────────
% BEAMER THEME & BASIC SETUP
% ─────────────────────────────────────────────────────────────────────────────
\usetheme{Madrid}
\usecolortheme{default}

% Remove navigation symbols
\beamertemplatenavigationsymbolsempty

% ─────────────────────────────────────────────────────────────────────────────
% COLORS — HIT Branding
% ─────────────────────────────────────────────────────────────────────────────
\definecolor{hitblue}{RGB}{0,51,102}        % Primary: HIT Navy
\definecolor{hitgold}{RGB}{192,148,0}       % Accent: HIT Gold
\definecolor{hitlight}{RGB}{230,240,250}    % Light background

% Override Beamer palette
\setbeamercolor{primary}{fg=white,bg=hitblue}
\setbeamercolor{secondary}{fg=hitblue,bg=hitlight}
\setbeamercolor{structure}{fg=hitblue}
\setbeamercolor{alerted text}{fg=hitgold}

% ─────────────────────────────────────────────────────────────────────────────
% PACKAGES
% ─────────────────────────────────────────────────────────────────────────────
\usepackage{amsmath}
\usepackage{amssymb}
\usepackage{mathtools}
\usepackage{booktabs}
\usepackage{graphicx}
\usepackage{hyperref}
\usepackage{tikz}
\usepackage{pgfplots}
\usepackage{tcolorbox}
\usepackage{listings}
\usepackage{xcolor}
\usepackage{algorithm}
\usepackage{algpseudocode}

% ─────────────────────────────────────────────────────────────────────────────
% TikZ LIBRARIES
% ─────────────────────────────────────────────────────────────────────────────
\usetikzlibrary{shapes,arrows,positioning,calc,chains,scopes}
\pgfplotsset{compat=1.17}

% ─────────────────────────────────────────────────────────────────────────────
% CODE LISTINGS
% ─────────────────────────────────────────────────────────────────────────────
\lstset{
  basicstyle=\ttfamily\small,
  breaklines=true,
  commentstyle=\color{gray},
  keywordstyle=\color{hitblue},
  stringstyle=\color{hitgold},
  showstringspaces=false,
  backgroundcolor=\color{hitlight},
  frame=single,
  rulecolor=\color{hitblue}
}

% ─────────────────────────────────────────────────────────────────────────────
% TCOLORBOX STYLES
% ─────────────────────────────────────────────────────────────────────────────
\tcbset{
  hitbox/.style={
    colback=hitlight,
    colframe=hitblue,
    fonttitle=\bfseries,
    boxrule=1pt
  },
  alertbox/.style={
    colback=yellow!10,
    colframe=hitgold,
    fonttitle=\bfseries,
    boxrule=1.5pt
  }
}

% ─────────────────────────────────────────────────────────────────────────────
% CUSTOM COMMANDS
% ─────────────────────────────────────────────────────────────────────────────

% Placeholder for figures
\newcommand{\placeholder}[3]{
  \begin{center}
    \fbox{\begin{minipage}[c][#2][c]{#1}
      \centering\small\color{gray}[Figure: #3]
    \end{minipage}}
  \end{center}
}

% Section divider frame
\newcommand{\sectionframe}[2]{
  \begin{frame}[plain]
    \begin{tikzpicture}[remember picture,overlay]
      \fill[hitblue] (current page.south west) rectangle (current page.north east);
      \node[anchor=center, white, font=\Large\bfseries] at (current page.center) {
        \begin{tabular}{c}
          #1 \\[0.3cm]
          {\small\color{hitgold}#2}
        \end{tabular}
      };
    \end{tikzpicture}
  \end{frame}
}

% ─────────────────────────────────────────────────────────────────────────────
% LOAD CUSTOM MACROS
% ─────────────────────────────────────────────────────────────────────────────
% ═══════════════════════════════════════════════════════════════════════════════
% SHARED MACROS — Mathematical Notation for Foundation Models Course
% ═══════════════════════════════════════════════════════════════════════════════

% ─────────────────────────────────────────────────────────────────────────────
% ACTIVATION & LAYER FUNCTIONS
% ─────────────────────────────────────────────────────────────────────────────
\newcommand{\softmax}{\mathrm{softmax}}
\newcommand{\sigmoid}{\mathrm{sigmoid}}
\newcommand{\relu}{\mathrm{ReLU}}
\newcommand{\gelu}{\mathrm{GELU}}
\newcommand{\LayerNorm}{\mathrm{LayerNorm}}
\newcommand{\FFN}{\mathrm{FFN}}
\newcommand{\MHA}{\mathrm{MHA}}
\newcommand{\Attn}{\mathrm{Attention}}

% ─────────────────────────────────────────────────────────────────────────────
% ATTENTION COMPONENTS
% ─────────────────────────────────────────────────────────────────────────────
\newcommand{\query}{Q}
\newcommand{\key}{K}
\newcommand{\val}{V}
\newcommand{\dmodel}{d_{\text{model}}}
\newcommand{\dk}{d_k}
\newcommand{\nheads}{h}

% Scaled dot-product attention formula
\newcommand{\SDPA}{
  \Attn(\query, \key, \val) = \softmax\left(\frac{\query \key^T}{\sqrt{\dk}}\right) \val
}

\newcommand{\CrossAttn}{\text{CrossAttention}}

% ─────────────────────────────────────────────────────────────────────────────
% PROBABILITY & EXPECTATION
% ─────────────────────────────────────────────────────────────────────────────
\newcommand{\E}{\mathbb{E}}
\newcommand{\Prob}{\mathbb{P}}
\newcommand{\KL}{\mathrm{KL}}
\newcommand{\ELBO}{\mathrm{ELBO}}
\newcommand{\given}{\mid}

% ─────────────────────────────────────────────────────────────────────────────
% NORMS & OPERATIONS
% ─────────────────────────────────────────────────────────────────────────────
\newcommand{\norm}[1]{\left\lVert #1 \right\rVert}
\newcommand{\abs}[1]{\left| #1 \right|}
\newcommand{\inner}[2]{\langle #1, #2 \rangle}
\newcommand{\T}{^{\top}}

% ─────────────────────────────────────────────────────────────────────────────
% VECTORS & MATRICES
% ─────────────────────────────────────────────────────────────────────────────
\newcommand{\bvec}[1]{\mathbf{#1}}
\newcommand{\mat}[1]{\mathbf{#1}}
\newcommand{\iid}{\stackrel{\text{i.i.d.}}{\sim}}
\newcommand{\defeq}{\coloneq}

% ─────────────────────────────────────────────────────────────────────────────
% LANGUAGE MODELING
% ─────────────────────────────────────────────────────────────────────────────
\newcommand{\vocab}{\mathcal{V}}
\newcommand{\context}{\text{context}}
\newcommand{\token}{t}
\newcommand{\seq}{x}
\newcommand{\ptheta}{p_\theta}
\newcommand{\pref}{p_{\text{ref}}}

% ─────────────────────────────────────────────────────────────────────────────
% RLHF & DPO
% ─────────────────────────────────────────────────────────────────────────────
\newcommand{\piref}{\pi_{\text{ref}}}
\newcommand{\pitheta}{\pi_\theta}
\newcommand{\yw}{y_w}
\newcommand{\yl}{y_l}
\newcommand{\DPOloss}{\mathcal{L}_{\mathrm{DPO}}}

% ─────────────────────────────────────────────────────────────────────────────
% RETRIEVAL & RAG
% ─────────────────────────────────────────────────────────────────────────────
\newcommand{\docs}{\mathcal{D}}
\newcommand{\qvec}{q}
\newcommand{\dvec}{d}

% ─────────────────────────────────────────────────────────────────────────────
% AGENTS & REASONING
% ─────────────────────────────────────────────────────────────────────────────
\newcommand{\obs}{o}
\newcommand{\act}{a}
\newcommand{\thought}{t}
\newcommand{\tools}{\mathcal{T}}
\newcommand{\history}{h}
\newcommand{\concat}{\oplus}

% ─────────────────────────────────────────────────────────────────────────────
% SETS & LOGIC
% ─────────────────────────────────────────────────────────────────────────────
\newcommand{\R}{\mathbb{R}}
\newcommand{\N}{\mathbb{N}}
\newcommand{\Z}{\mathbb{Z}}

\endinput


% ─────────────────────────────────────────────────────────────────────────────
% TITLE & AUTHOR (can be overridden in individual lectures)
% ─────────────────────────────────────────────────────────────────────────────
\author[D. Lederman]{Dror Lederman}
\institute{Holon Institute of Technology (HIT) \\ Data Science Department}
\date{\today}

% ─────────────────────────────────────────────────────────────────────────────
% FOOTER & HEADER
% ─────────────────────────────────────────────────────────────────────────────
\setbeamertemplate{footline}{
  \leavevmode%
  \hbox{%
    \begin{beamercolorbox}[wd=0.33\paperwidth,ht=2.25ex,dp=1ex,center]{author in head/foot}%
      \usebeamerfont{author in head/foot}\insertshortauthor
    \end{beamercolorbox}%
    \begin{beamercolorbox}[wd=0.34\paperwidth,ht=2.25ex,dp=1ex,center]{title in head/foot}%
      \usebeamerfont{title in head/foot}\insertshorttitle
    \end{beamercolorbox}%
    \begin{beamercolorbox}[wd=0.33\paperwidth,ht=2.25ex,dp=1ex,right]{frame number}%
      \usebeamerfont{frame number}\insertframenumber{} / \inserttotalframenumber\hspace*{2ex}
    \end{beamercolorbox}%
  }%
  \vskip0pt%
}

\endinput


\title{Week 2: Transformer Architectures Revisited}
\subtitle{Foundation Models and Agentic AI}
\author{Lecturer Name}
\institute{Holon Institute of Technology (HIT)}
\date{Semester, 2025}

\begin{document}

\begin{frame}
  \titlepage
\end{frame}

\begin{frame}{Learning Objectives}
  After this lecture you will be able to:
  \begin{itemize}
    \item Derive the \textbf{scaled dot-product attention} and multi-head attention from first principles.
    \item Compare and contrast \textbf{absolute, rotary (RoPE), and ALiBi} positional encodings.
    \item Explain the design choices behind \textbf{decoder-only} vs.\ encoder-decoder architectures.
    \item Describe \textbf{sparse attention} patterns and when they are advantageous.
    \item Understand \textbf{Mixture of Experts (MoE)} and how it scales model capacity efficiently.
  \end{itemize}
\end{frame}

\section{Self-Attention from First Principles}
\sectionframe{Self-Attention from First Principles}{The core operation of every modern LLM}

\begin{frame}{Attention as Soft Retrieval}
  \begin{columns}[T]
    \column{0.52\textwidth}
      Think of attention as a \textbf{differentiable key-value store}:
      \begin{itemize}
        \item A \textbf{query} $\query_i$ represents: ``what am I looking for?''
        \item \textbf{Keys} $\key_j$ represent: ``what does token $j$ offer?''
        \item \textbf{Values} $\val_j$ represent: ``what information does token $j$ carry?''
      \end{itemize}
      \vspace{0.4em}
      Retrieval is \emph{soft}: all values are aggregated,
      weighted by query-key similarity.
    \column{0.44\textwidth}
      % tikz/attention_head.tex
% Scaled dot-product attention diagram
% Usage: % tikz/attention_head.tex
% Scaled dot-product attention diagram
% Usage: % tikz/attention_head.tex
% Scaled dot-product attention diagram
% Usage: \input{../../tikz/attention_head}
\begin{tikzpicture}[
  node distance=0.7cm and 1.0cm,
  box/.style={rectangle, draw=hitblue, fill=hitlightblue, rounded corners=2pt,
              minimum width=1.6cm, minimum height=0.55cm, font=\small, align=center},
  op/.style={circle, draw=hitgold, fill=hitgold!20,
             minimum size=0.55cm, font=\small\bfseries},
  arr/.style={->, >=Stealth, thick, hitblue},
]

% Input nodes
\node[box] (Q) {$\query$};
\node[box, right=1.6cm of Q] (K) {$\key$};
\node[box, right=1.6cm of K] (V) {$\val$};

% MatMul QK^T
\node[op, below=0.8cm of $(Q)!0.5!(K)$] (matmul1) {$\times$};

% Scale
\node[box, below=0.7cm of matmul1] (scale) {Scale $\frac{1}{\sqrt{\dk}}$};

% Mask (optional)
\node[box, below=0.7cm of scale, dashed, draw=gray, fill=gray!8, font=\small\itshape]
  (mask) {Mask (opt.)};

% Softmax
\node[box, below=0.7cm of mask] (sm) {Softmax};

% MatMul with V
\node[op, below=0.8cm of $(sm)!0.5!(V)+(1.6cm,0)$] (matmul2) {$\times$};

% Output
\node[box, below=0.7cm of matmul2] (out) {\textbf{Output}};

% Arrows
\draw[arr] (Q.south) -- ++(0,-0.3) -| (matmul1.north west);
\draw[arr] (K.south) -- ++(0,-0.3) -| (matmul1.north east);
\draw[arr] (matmul1) -- (scale);
\draw[arr] (scale) -- (mask);
\draw[arr] (mask) -- (sm);
\draw[arr] (sm.south) -- ++(0,-0.3) -| (matmul2.north west);
\draw[arr] (V.south) -- ++(0,-2.9) -| (matmul2.north east);
\draw[arr] (matmul2) -- (out);

% Label
\node[font=\footnotesize\itshape, gray, above=0.05cm of Q] {Queries};
\node[font=\footnotesize\itshape, gray, above=0.05cm of K] {Keys};
\node[font=\footnotesize\itshape, gray, above=0.05cm of V] {Values};

\end{tikzpicture}

\begin{tikzpicture}[
  node distance=0.7cm and 1.0cm,
  box/.style={rectangle, draw=hitblue, fill=hitlightblue, rounded corners=2pt,
              minimum width=1.6cm, minimum height=0.55cm, font=\small, align=center},
  op/.style={circle, draw=hitgold, fill=hitgold!20,
             minimum size=0.55cm, font=\small\bfseries},
  arr/.style={->, >=Stealth, thick, hitblue},
]

% Input nodes
\node[box] (Q) {$\query$};
\node[box, right=1.6cm of Q] (K) {$\key$};
\node[box, right=1.6cm of K] (V) {$\val$};

% MatMul QK^T
\node[op, below=0.8cm of $(Q)!0.5!(K)$] (matmul1) {$\times$};

% Scale
\node[box, below=0.7cm of matmul1] (scale) {Scale $\frac{1}{\sqrt{\dk}}$};

% Mask (optional)
\node[box, below=0.7cm of scale, dashed, draw=gray, fill=gray!8, font=\small\itshape]
  (mask) {Mask (opt.)};

% Softmax
\node[box, below=0.7cm of mask] (sm) {Softmax};

% MatMul with V
\node[op, below=0.8cm of $(sm)!0.5!(V)+(1.6cm,0)$] (matmul2) {$\times$};

% Output
\node[box, below=0.7cm of matmul2] (out) {\textbf{Output}};

% Arrows
\draw[arr] (Q.south) -- ++(0,-0.3) -| (matmul1.north west);
\draw[arr] (K.south) -- ++(0,-0.3) -| (matmul1.north east);
\draw[arr] (matmul1) -- (scale);
\draw[arr] (scale) -- (mask);
\draw[arr] (mask) -- (sm);
\draw[arr] (sm.south) -- ++(0,-0.3) -| (matmul2.north west);
\draw[arr] (V.south) -- ++(0,-2.9) -| (matmul2.north east);
\draw[arr] (matmul2) -- (out);

% Label
\node[font=\footnotesize\itshape, gray, above=0.05cm of Q] {Queries};
\node[font=\footnotesize\itshape, gray, above=0.05cm of K] {Keys};
\node[font=\footnotesize\itshape, gray, above=0.05cm of V] {Values};

\end{tikzpicture}

\begin{tikzpicture}[
  node distance=0.7cm and 1.0cm,
  box/.style={rectangle, draw=hitblue, fill=hitlightblue, rounded corners=2pt,
              minimum width=1.6cm, minimum height=0.55cm, font=\small, align=center},
  op/.style={circle, draw=hitgold, fill=hitgold!20,
             minimum size=0.55cm, font=\small\bfseries},
  arr/.style={->, >=Stealth, thick, hitblue},
]

% Input nodes
\node[box] (Q) {$\query$};
\node[box, right=1.6cm of Q] (K) {$\key$};
\node[box, right=1.6cm of K] (V) {$\val$};

% MatMul QK^T
\node[op, below=0.8cm of $(Q)!0.5!(K)$] (matmul1) {$\times$};

% Scale
\node[box, below=0.7cm of matmul1] (scale) {Scale $\frac{1}{\sqrt{\dk}}$};

% Mask (optional)
\node[box, below=0.7cm of scale, dashed, draw=gray, fill=gray!8, font=\small\itshape]
  (mask) {Mask (opt.)};

% Softmax
\node[box, below=0.7cm of mask] (sm) {Softmax};

% MatMul with V
\node[op, below=0.8cm of $(sm)!0.5!(V)+(1.6cm,0)$] (matmul2) {$\times$};

% Output
\node[box, below=0.7cm of matmul2] (out) {\textbf{Output}};

% Arrows
\draw[arr] (Q.south) -- ++(0,-0.3) -| (matmul1.north west);
\draw[arr] (K.south) -- ++(0,-0.3) -| (matmul1.north east);
\draw[arr] (matmul1) -- (scale);
\draw[arr] (scale) -- (mask);
\draw[arr] (mask) -- (sm);
\draw[arr] (sm.south) -- ++(0,-0.3) -| (matmul2.north west);
\draw[arr] (V.south) -- ++(0,-2.9) -| (matmul2.north east);
\draw[arr] (matmul2) -- (out);

% Label
\node[font=\footnotesize\itshape, gray, above=0.05cm of Q] {Queries};
\node[font=\footnotesize\itshape, gray, above=0.05cm of K] {Keys};
\node[font=\footnotesize\itshape, gray, above=0.05cm of V] {Values};

\end{tikzpicture}

  \end{columns}
  \note{The key-value store analogy is pedagogically useful. In a database, you retrieve one exact match. In attention, you retrieve a weighted blend.}
\end{frame}

\begin{frame}{Scaled Dot-Product Attention}
  Given input sequence $\mat{X} \in \R^{n \times \dmodel}$:
  \[
    \query = \mat{X}\mat{W}^Q, \quad
    \key   = \mat{X}\mat{W}^K, \quad
    \val   = \mat{X}\mat{W}^V
  \]
  \vspace{-0.3em}
  \begin{tcolorbox}[hitbox, title=Scaled Dot-Product Attention]
    \[
      \Attn(\query, \key, \val)
      = \softmax\!\left(\frac{\query\key\T}{\sqrt{\dk}}\right)\val
    \]
  \end{tcolorbox}
  \vspace{0.3em}
  \textbf{Why scale by $\sqrt{\dk}$?}\\
  Dot products grow in magnitude with $\dk$ (variance $= \dk$).
  Without scaling, softmax saturates $\Rightarrow$ vanishing gradients.
  \note{Show this derivation step by step. Students from the DL course know softmax, but the scaling argument often needs re-emphasis.}
\end{frame}

\begin{frame}{Multi-Head Attention}
  Run $h$ attention heads \emph{in parallel} with different projections:
  \[
    \text{head}_i = \Attn\!\left(\query\mat{W}_i^Q,\ \key\mat{W}_i^K,\ \val\mat{W}_i^V\right)
  \]
  \[
    \MHA(\query,\key,\val) = \left[\text{head}_1;\ldots;\text{head}_h\right]\mat{W}^O
  \]
  \vspace{0.4em}
  \textbf{Why multiple heads?}
  \begin{itemize}
    \item Each head can attend to different positional/semantic relations simultaneously.
    \item Head diversity is \emph{empirically} important — pruning heads degrades performance.
    \item Typical: $h = 16$ heads, $\dk = \dmodel / h = 64$.
  \end{itemize}
  \note{In practice, all heads are computed in one batched matrix multiplication. The diagram shows them separately for clarity.}
\end{frame}

\begin{frame}{The Full Transformer Block}
  \begin{columns}[T]
    \column{0.50\textwidth}
      \textbf{One transformer layer:}
      \begin{align*}
        \tilde{\mat{X}} &= \LayerNorm\!\left(\mat{X} + \MHA(\mat{X})\right)\\
        \mat{X}' &= \LayerNorm\!\left(\tilde{\mat{X}} + \FFN(\tilde{\mat{X}})\right)
      \end{align*}
      \vspace{0.3em}
      Where $\FFN(\bvec{x}) = \max(0,\bvec{x}\mat{W}_1 + \bvec{b}_1)\mat{W}_2 + \bvec{b}_2$
      \vspace{0.5em}\\
      \textbf{Pre-norm vs.\ post-norm:}
      \begin{itemize}
        \item Original Vaswani: post-norm
        \item GPT/Llama family: pre-norm (more stable training)
      \end{itemize}
    \column{0.46\textwidth}
      % tikz/transformer_arch.tex
% Transformer encoder-decoder architecture (simplified)
% Usage: % tikz/transformer_arch.tex
% Transformer encoder-decoder architecture (simplified)
% Usage: % tikz/transformer_arch.tex
% Transformer encoder-decoder architecture (simplified)
% Usage: \input{../../tikz/transformer_arch}
\begin{tikzpicture}[
  node distance=0.45cm and 1.8cm,
  block/.style={rectangle, draw=hitblue, fill=hitlightblue, rounded corners=3pt,
                minimum width=2.4cm, minimum height=0.55cm, font=\footnotesize, align=center},
  sublayer/.style={rectangle, draw=hitblue!60, fill=hitblue!8, rounded corners=2pt,
                   minimum width=2.4cm, minimum height=0.5cm, font=\scriptsize, align=center},
  io/.style={rectangle, draw=hitgold, fill=hitgold!15, rounded corners=3pt,
             minimum width=2.4cm, minimum height=0.5cm, font=\footnotesize, align=center},
  arr/.style={->, >=Stealth, thick, hitblue},
  rarr/.style={->, >=Stealth, thick, hitgold},
  label/.style={font=\scriptsize\bfseries, hitblue},
]

%% ── Encoder (left column) ────────────────────────────────────────────────────
\node[io] (enc_in) {Input Embeddings};
\node[block, above=0.5cm of enc_in] (enc_pe) {+ Positional Encoding};

% Encoder block (repeated Nx)
\node[sublayer, above=0.6cm of enc_pe] (enc_mha) {Multi-Head\\Self-Attention};
\node[sublayer, above=0.45cm of enc_mha] (enc_add1) {Add \& LayerNorm};
\node[sublayer, above=0.45cm of enc_add1] (enc_ffn) {Feed-Forward};
\node[sublayer, above=0.45cm of enc_ffn] (enc_add2) {Add \& LayerNorm};

% Encoder Nx box
\node[draw=hitblue, dashed, fit=(enc_mha)(enc_add2),
      inner sep=4pt, label={[label]right:$N\times$}] (encbox) {};

\node[io, above=0.5cm of enc_add2] (enc_out) {Encoder Output};

% Encoder arrows
\draw[arr] (enc_in) -- (enc_pe);
\draw[arr] (enc_pe) -- (enc_mha);
\draw[arr] (enc_mha) -- (enc_add1);
\draw[arr] (enc_add1) -- (enc_ffn);
\draw[arr] (enc_ffn) -- (enc_add2);
\draw[arr] (enc_add2) -- (enc_out);

%% ── Decoder (right column) ───────────────────────────────────────────────────
\node[io, right=2.0cm of enc_in] (dec_in) {Output Embeddings};
\node[block, above=0.5cm of dec_in] (dec_pe) {+ Positional Encoding};

% Decoder block
\node[sublayer, above=0.6cm of dec_pe] (dec_mmha) {Masked MH\\Self-Attention};
\node[sublayer, above=0.45cm of dec_mmha] (dec_add1) {Add \& LayerNorm};
\node[sublayer, above=0.45cm of dec_add1] (dec_cross) {Cross-Attention\\(Enc $\to$ Dec)};
\node[sublayer, above=0.45cm of dec_cross] (dec_add2) {Add \& LayerNorm};
\node[sublayer, above=0.45cm of dec_add2] (dec_ffn) {Feed-Forward};
\node[sublayer, above=0.45cm of dec_ffn] (dec_add3) {Add \& LayerNorm};

\node[draw=hitblue, dashed, fit=(dec_mmha)(dec_add3),
      inner sep=4pt, label={[label]right:$N\times$}] (decbox) {};

\node[io, above=0.5cm of dec_add3] (dec_out) {Linear + Softmax};
\node[io, above=0.45cm of dec_out] (dec_prob) {Output Probs};

% Decoder arrows
\draw[arr] (dec_in) -- (dec_pe);
\draw[arr] (dec_pe) -- (dec_mmha);
\draw[arr] (dec_mmha) -- (dec_add1);
\draw[arr] (dec_add1) -- (dec_cross);
\draw[arr] (dec_cross) -- (dec_add2);
\draw[arr] (dec_add2) -- (dec_ffn);
\draw[arr] (dec_ffn) -- (dec_add3);
\draw[arr] (dec_add3) -- (dec_out);
\draw[arr] (dec_out) -- (dec_prob);

% Cross-attention arrow from encoder
\draw[rarr, dashed] (enc_out.east) -- ++(0.4,0) |-
  node[above, font=\scriptsize, hitgold, pos=0.7] {K, V} (dec_cross.west);

%% ── Column labels ────────────────────────────────────────────────────────────
\node[font=\small\bfseries, hitblue, above=0.1cm of dec_prob] {};
\node[font=\footnotesize, gray, below=0.2cm of enc_in] {\textbf{Encoder}};
\node[font=\footnotesize, gray, below=0.2cm of dec_in] {\textbf{Decoder}};

\end{tikzpicture}

\begin{tikzpicture}[
  node distance=0.45cm and 1.8cm,
  block/.style={rectangle, draw=hitblue, fill=hitlightblue, rounded corners=3pt,
                minimum width=2.4cm, minimum height=0.55cm, font=\footnotesize, align=center},
  sublayer/.style={rectangle, draw=hitblue!60, fill=hitblue!8, rounded corners=2pt,
                   minimum width=2.4cm, minimum height=0.5cm, font=\scriptsize, align=center},
  io/.style={rectangle, draw=hitgold, fill=hitgold!15, rounded corners=3pt,
             minimum width=2.4cm, minimum height=0.5cm, font=\footnotesize, align=center},
  arr/.style={->, >=Stealth, thick, hitblue},
  rarr/.style={->, >=Stealth, thick, hitgold},
  label/.style={font=\scriptsize\bfseries, hitblue},
]

%% ── Encoder (left column) ────────────────────────────────────────────────────
\node[io] (enc_in) {Input Embeddings};
\node[block, above=0.5cm of enc_in] (enc_pe) {+ Positional Encoding};

% Encoder block (repeated Nx)
\node[sublayer, above=0.6cm of enc_pe] (enc_mha) {Multi-Head\\Self-Attention};
\node[sublayer, above=0.45cm of enc_mha] (enc_add1) {Add \& LayerNorm};
\node[sublayer, above=0.45cm of enc_add1] (enc_ffn) {Feed-Forward};
\node[sublayer, above=0.45cm of enc_ffn] (enc_add2) {Add \& LayerNorm};

% Encoder Nx box
\node[draw=hitblue, dashed, fit=(enc_mha)(enc_add2),
      inner sep=4pt, label={[label]right:$N\times$}] (encbox) {};

\node[io, above=0.5cm of enc_add2] (enc_out) {Encoder Output};

% Encoder arrows
\draw[arr] (enc_in) -- (enc_pe);
\draw[arr] (enc_pe) -- (enc_mha);
\draw[arr] (enc_mha) -- (enc_add1);
\draw[arr] (enc_add1) -- (enc_ffn);
\draw[arr] (enc_ffn) -- (enc_add2);
\draw[arr] (enc_add2) -- (enc_out);

%% ── Decoder (right column) ───────────────────────────────────────────────────
\node[io, right=2.0cm of enc_in] (dec_in) {Output Embeddings};
\node[block, above=0.5cm of dec_in] (dec_pe) {+ Positional Encoding};

% Decoder block
\node[sublayer, above=0.6cm of dec_pe] (dec_mmha) {Masked MH\\Self-Attention};
\node[sublayer, above=0.45cm of dec_mmha] (dec_add1) {Add \& LayerNorm};
\node[sublayer, above=0.45cm of dec_add1] (dec_cross) {Cross-Attention\\(Enc $\to$ Dec)};
\node[sublayer, above=0.45cm of dec_cross] (dec_add2) {Add \& LayerNorm};
\node[sublayer, above=0.45cm of dec_add2] (dec_ffn) {Feed-Forward};
\node[sublayer, above=0.45cm of dec_ffn] (dec_add3) {Add \& LayerNorm};

\node[draw=hitblue, dashed, fit=(dec_mmha)(dec_add3),
      inner sep=4pt, label={[label]right:$N\times$}] (decbox) {};

\node[io, above=0.5cm of dec_add3] (dec_out) {Linear + Softmax};
\node[io, above=0.45cm of dec_out] (dec_prob) {Output Probs};

% Decoder arrows
\draw[arr] (dec_in) -- (dec_pe);
\draw[arr] (dec_pe) -- (dec_mmha);
\draw[arr] (dec_mmha) -- (dec_add1);
\draw[arr] (dec_add1) -- (dec_cross);
\draw[arr] (dec_cross) -- (dec_add2);
\draw[arr] (dec_add2) -- (dec_ffn);
\draw[arr] (dec_ffn) -- (dec_add3);
\draw[arr] (dec_add3) -- (dec_out);
\draw[arr] (dec_out) -- (dec_prob);

% Cross-attention arrow from encoder
\draw[rarr, dashed] (enc_out.east) -- ++(0.4,0) |-
  node[above, font=\scriptsize, hitgold, pos=0.7] {K, V} (dec_cross.west);

%% ── Column labels ────────────────────────────────────────────────────────────
\node[font=\small\bfseries, hitblue, above=0.1cm of dec_prob] {};
\node[font=\footnotesize, gray, below=0.2cm of enc_in] {\textbf{Encoder}};
\node[font=\footnotesize, gray, below=0.2cm of dec_in] {\textbf{Decoder}};

\end{tikzpicture}

\begin{tikzpicture}[
  node distance=0.45cm and 1.8cm,
  block/.style={rectangle, draw=hitblue, fill=hitlightblue, rounded corners=3pt,
                minimum width=2.4cm, minimum height=0.55cm, font=\footnotesize, align=center},
  sublayer/.style={rectangle, draw=hitblue!60, fill=hitblue!8, rounded corners=2pt,
                   minimum width=2.4cm, minimum height=0.5cm, font=\scriptsize, align=center},
  io/.style={rectangle, draw=hitgold, fill=hitgold!15, rounded corners=3pt,
             minimum width=2.4cm, minimum height=0.5cm, font=\footnotesize, align=center},
  arr/.style={->, >=Stealth, thick, hitblue},
  rarr/.style={->, >=Stealth, thick, hitgold},
  label/.style={font=\scriptsize\bfseries, hitblue},
]

%% ── Encoder (left column) ────────────────────────────────────────────────────
\node[io] (enc_in) {Input Embeddings};
\node[block, above=0.5cm of enc_in] (enc_pe) {+ Positional Encoding};

% Encoder block (repeated Nx)
\node[sublayer, above=0.6cm of enc_pe] (enc_mha) {Multi-Head\\Self-Attention};
\node[sublayer, above=0.45cm of enc_mha] (enc_add1) {Add \& LayerNorm};
\node[sublayer, above=0.45cm of enc_add1] (enc_ffn) {Feed-Forward};
\node[sublayer, above=0.45cm of enc_ffn] (enc_add2) {Add \& LayerNorm};

% Encoder Nx box
\node[draw=hitblue, dashed, fit=(enc_mha)(enc_add2),
      inner sep=4pt, label={[label]right:$N\times$}] (encbox) {};

\node[io, above=0.5cm of enc_add2] (enc_out) {Encoder Output};

% Encoder arrows
\draw[arr] (enc_in) -- (enc_pe);
\draw[arr] (enc_pe) -- (enc_mha);
\draw[arr] (enc_mha) -- (enc_add1);
\draw[arr] (enc_add1) -- (enc_ffn);
\draw[arr] (enc_ffn) -- (enc_add2);
\draw[arr] (enc_add2) -- (enc_out);

%% ── Decoder (right column) ───────────────────────────────────────────────────
\node[io, right=2.0cm of enc_in] (dec_in) {Output Embeddings};
\node[block, above=0.5cm of dec_in] (dec_pe) {+ Positional Encoding};

% Decoder block
\node[sublayer, above=0.6cm of dec_pe] (dec_mmha) {Masked MH\\Self-Attention};
\node[sublayer, above=0.45cm of dec_mmha] (dec_add1) {Add \& LayerNorm};
\node[sublayer, above=0.45cm of dec_add1] (dec_cross) {Cross-Attention\\(Enc $\to$ Dec)};
\node[sublayer, above=0.45cm of dec_cross] (dec_add2) {Add \& LayerNorm};
\node[sublayer, above=0.45cm of dec_add2] (dec_ffn) {Feed-Forward};
\node[sublayer, above=0.45cm of dec_ffn] (dec_add3) {Add \& LayerNorm};

\node[draw=hitblue, dashed, fit=(dec_mmha)(dec_add3),
      inner sep=4pt, label={[label]right:$N\times$}] (decbox) {};

\node[io, above=0.5cm of dec_add3] (dec_out) {Linear + Softmax};
\node[io, above=0.45cm of dec_out] (dec_prob) {Output Probs};

% Decoder arrows
\draw[arr] (dec_in) -- (dec_pe);
\draw[arr] (dec_pe) -- (dec_mmha);
\draw[arr] (dec_mmha) -- (dec_add1);
\draw[arr] (dec_add1) -- (dec_cross);
\draw[arr] (dec_cross) -- (dec_add2);
\draw[arr] (dec_add2) -- (dec_ffn);
\draw[arr] (dec_ffn) -- (dec_add3);
\draw[arr] (dec_add3) -- (dec_out);
\draw[arr] (dec_out) -- (dec_prob);

% Cross-attention arrow from encoder
\draw[rarr, dashed] (enc_out.east) -- ++(0.4,0) |-
  node[above, font=\scriptsize, hitgold, pos=0.7] {K, V} (dec_cross.west);

%% ── Column labels ────────────────────────────────────────────────────────────
\node[font=\small\bfseries, hitblue, above=0.1cm of dec_prob] {};
\node[font=\footnotesize, gray, below=0.2cm of enc_in] {\textbf{Encoder}};
\node[font=\footnotesize, gray, below=0.2cm of dec_in] {\textbf{Decoder}};

\end{tikzpicture}

  \end{columns}
  \note{Emphasize residual connections: they are critical for training deep networks. Without them, gradients vanish.}
\end{frame}

\section{Positional Encoding}
\sectionframe{Positional Encoding}{Giving the model a sense of position}

\begin{frame}{Why Positional Encoding?}
  \begin{tcolorbox}[alertbox, title=Problem]
    Attention is \textbf{permutation-invariant}: shuffling tokens produces the same attention
    weights. We must inject position information.
  \end{tcolorbox}
  \vspace{0.5em}
  \textbf{Three main approaches:}
  \begin{enumerate}
    \item \textbf{Absolute (sinusoidal):} Fixed frequencies, added to embeddings. (Vaswani 2017)
    \item \textbf{Learned absolute:} Learn a lookup table of position embeddings. (BERT, GPT-2)
    \item \textbf{Relative:} Modify attention scores based on relative distance. (T5, RoPE, ALiBi)
  \end{enumerate}
\end{frame}

\begin{frame}{Rotary Position Embedding (RoPE)}
  \textbf{Su et al.\ (2024)}~\cite{su2024roformer} — used in Llama, Mistral, GPT-NeoX.
  \vspace{0.4em}
  \begin{tcolorbox}[hitbox, title=Core Idea]
    Rotate the query and key vectors by an angle proportional to their absolute position.
    The dot product $\query_m \cdot \key_n$ then depends only on \emph{relative distance} $m - n$.
  \end{tcolorbox}
  \vspace{0.4em}
  For a 2D vector at position $m$, dimension pair $(2i, 2i{+}1)$:
  \[
    f(\bvec{x}, m) = \begin{pmatrix} x_{2i}\cos(m\theta_i) - x_{2i+1}\sin(m\theta_i)\\
                                      x_{2i}\sin(m\theta_i) + x_{2i+1}\cos(m\theta_i) \end{pmatrix}
  \]
  \textbf{Advantages:} Stable at long contexts; extrapolates better than absolute PE.
  \note{Students with linear algebra background will appreciate this as a rotation matrix. The key property: dot product of rotated vectors = relative-distance-dependent similarity.}
\end{frame}

\begin{frame}{ALiBi: Linear Biases for Extrapolation}
  \textbf{Press et al.\ (2022)}~\cite{press2022alibi} — used in MPT, BloomZ.
  \vspace{0.4em}
  \begin{tcolorbox}[hitbox, title=ALiBi Mechanism]
    Instead of adding positional embeddings, \textbf{subtract} a linear penalty from attention
    scores proportional to relative distance:
    \[
      A_{ij} = \query_i \cdot \key_j\T / \sqrt{\dk} - m_h \cdot |i - j|
    \]
    Slope $m_h$ varies per head (geometric sequence).
  \end{tcolorbox}
  \vspace{0.3em}
  \textbf{Key property:} Train on short contexts, test on much longer ones — \emph{length extrapolation}.
\end{frame}

\section{Decoder-Only vs.\ Encoder-Decoder}
\sectionframe{Architecture Families}{Encoder-only, decoder-only, encoder-decoder}

\begin{frame}{Architecture Comparison}
  \begin{center}
  \begin{tabular}{llll}
    \toprule
    \textbf{Architecture} & \textbf{Attention mask} & \textbf{Best for} & \textbf{Examples} \\
    \midrule
    Encoder-only & Bidirectional & Classification, NER & BERT, RoBERTa \\
    Decoder-only & Causal (left-to-right) & Generation, in-context & GPT, Llama, Claude \\
    Encoder-decoder & Mixed & Seq2seq, summarization & T5, BART, Gemini \\
    \bottomrule
  \end{tabular}
  \end{center}
  \vspace{0.5em}
  \textbf{The trend:} Modern foundation models are almost exclusively \textbf{decoder-only}.\\
  \emph{Why?} Simpler architecture, scales better with next-token pretraining, supports in-context learning.
  \note{BERT is still widely used for embedding/retrieval tasks. But GPT-style models dominate generation.}
\end{frame}

\section{Sparse Attention \& Long Context}
\sectionframe{Sparse Attention \& Long Contexts}{Scaling beyond the quadratic bottleneck}

\begin{frame}{The Quadratic Problem}
  \begin{tcolorbox}[alertbox, title=Computational Bottleneck]
    Full self-attention computes $\query\key\T \in \R^{n \times n}$. \\
    Time: $O(n^2 d)$ \quad Memory: $O(n^2)$ \\
    At $n = 128{,}000$ tokens: \textasciitilde 16 GB just for the attention matrix.
  \end{tcolorbox}
  \vspace{0.5em}
  \textbf{Solutions:}
  \begin{itemize}
    \item \textbf{Sparse attention} (Child et al.\ 2019~\cite{child2019sparsetransformer}): attend to a subset of tokens (strided, local, global)
    \item \textbf{Linear attention} (Katharopoulos et al.\ 2020): kernel approximation, $O(n)$
    \item \textbf{FlashAttention} (Dao et al.\ 2022): IO-aware exact attention, GPU-efficient
    \item \textbf{Sliding window / Ring attention}: attend to local windows + global tokens
  \end{itemize}
\end{frame}

\section{Mixture of Experts}
\sectionframe{Mixture of Experts}{Scaling capacity without scaling compute}

\begin{frame}{Mixture of Experts (MoE)}
  \textbf{Shazeer et al.\ (2017)}~\cite{shazeer2017moe}; revived in Mixtral~\cite{jiang2024mixtral}.
  \vspace{0.4em}
  \begin{tcolorbox}[hitbox, title=MoE Layer]
    Replace the FFN in each transformer block with $E$ expert networks $\{F_1,\ldots,F_E\}$.
    A lightweight \textbf{router} $G(\bvec{x})$ selects the top-$k$ experts per token:
    \[
      \text{MoE}(\bvec{x}) = \sum_{i \in \text{Top-}k} G(\bvec{x})_i \cdot F_i(\bvec{x})
    \]
  \end{tcolorbox}
  \vspace{0.4em}
  \textbf{Benefits:}
  \begin{itemize}
    \item Large \emph{total} parameter count, but only a fraction activated per token.
    \item Mixtral 8×7B: 47B total params, 13B active per forward pass.
    \item Comparable quality to dense 70B at lower inference cost.
  \end{itemize}
  \note{Challenge: load balancing. Without auxiliary loss, the router collapses to always selecting the same experts.}
\end{frame}

\begin{frame}{MoE: Load Balancing and Challenges}
  \begin{columns}[T]
    \column{0.55\textwidth}
      \textbf{Training challenges:}
      \begin{itemize}
        \item \textbf{Load imbalance}: some experts always selected.
          Fix: auxiliary load-balancing loss.
        \item \textbf{Routing instability}: router changes during training.
        \item \textbf{Communication overhead}: in distributed training, routing across devices.
      \end{itemize}
      \vspace{0.4em}
      \textbf{Inference challenges:}
      \begin{itemize}
        \item All experts must be in memory (even if only 2 are used).
        \item Batching efficiency drops when tokens route differently.
      \end{itemize}
    \column{0.42\textwidth}
      \placeholder{0.95\textwidth}{3.5cm}{MoE routing: 8 experts, top-2 selected per token}
  \end{columns}
\end{frame}

\section{Paper Discussion}
\begin{frame}{Paper Discussion}
  \papercite{Attention Is All You Need}
    {Vaswani, Shazeer, Parmar, Uszkoreit et al.}
    {NeurIPS 2017~\cite{vaswani2017attention}}
  \vspace{0.4em}
  \textbf{What it introduced:}
  \begin{itemize}
    \item The transformer architecture (encoder + decoder, multi-head attention, FFN, residual+norm)
    \item State-of-the-art on WMT translation with less compute than RNNs
    \item Position-wise FFN as a learned ``memory'' per position
  \end{itemize}
  \vspace{0.3em}
  \textbf{What it didn't predict:}
  \begin{itemize}
    \item Decoder-only models would dominate at scale
    \item The architecture would require almost no changes to scale to 100B+ parameters
  \end{itemize}
\end{frame}

\section{Lab / Demo}
\begin{frame}{Lab Idea: Attention Visualization}
  \begin{tcolorbox}[hitbox, title=Lab 2 — Attention Pattern Analysis]
    \textbf{Goal:} Understand what different attention heads learn.\\[4pt]
    \textbf{Tools:} \texttt{BertViz}, HuggingFace \texttt{transformers}
    \begin{enumerate}
      \item Load a pre-trained BERT-base model.
      \item Visualize attention heads on several sentences.
      \item Identify heads that attend to: (a) syntactic relations, (b) next/previous token, (c) special tokens (\texttt{[CLS]}, \texttt{[SEP]}).
      \item Repeat for GPT-2 (causal): do patterns differ?
    \end{enumerate}
    \textbf{Deliverable:} Annotated attention maps with 3-sentence interpretation per pattern.
  \end{tcolorbox}
\end{frame}

\section{Discussion \& Summary}
\begin{frame}{Discussion Questions}
  \begin{enumerate}
    \item Self-attention has $O(n^2)$ complexity. For which real-world applications is this the binding constraint? What are the engineering trade-offs of sparse vs.\ flash attention?
    \item Why did decoder-only models ``win'' despite encoder-decoder architectures existing for longer? Is there any task where encoder-only or encoder-decoder is still preferable?
    \item MoE models have more parameters but lower FLOPs. Does this mean they are ``bigger'' or ``smaller'' than a dense model of the same active parameter count?
    \item How does the choice of positional encoding affect a model's ability to generalize to longer sequences at inference time?
  \end{enumerate}
\end{frame}

\begin{frame}{Summary}
  \begin{itemize}
    \item \textbf{Self-attention} is a soft, differentiable key-value lookup; scaling requires division by $\sqrt{\dk}$.
    \item \textbf{Multi-head attention} allows simultaneous focus on different relational aspects.
    \item \textbf{Positional encoding} is necessary; modern approaches (RoPE, ALiBi) extrapolate better than sinusoidal.
    \item \textbf{Decoder-only} architectures dominate large-scale generative models.
    \item \textbf{Sparse attention} and \textbf{FlashAttention} address the $O(n^2)$ bottleneck.
    \item \textbf{MoE} decouples model capacity from per-token compute cost.
  \end{itemize}
  \vspace{0.3em}
  \textbf{Next week:} How we \emph{train} these models — pretraining objectives, RLHF, and DPO.
\end{frame}

\begin{frame}[allowframebreaks]{Suggested Reading}
  \nocite{vaswani2017attention, su2024roformer, press2022alibi,
          child2019sparsetransformer, shazeer2017moe, jiang2024mixtral}
  \bibliography{../../references}
\end{frame}

\end{document}
